\begin{table}[H]
\centering\small
\caption{Within- and between-party relation of the two scores of Figure~\ref{fig:party_map}: mean preferred policy for alcohol and tobacco regressed on mean preferred policy for the illicit drugs (both 1 = ban \ldots{} 9 = free sale), split into the party mean (between) and the voter's distance from it (within). Linear mixed models with a random intercept per party, survey weights; estimate (standard error), $p$ with Satterthwaite degrees of freedom. Between $-$ within: the party-mean coefficient in a model with the voter's own score, which tests whether the two slopes differ. Six largest parties: GroenLinks-PvdA, D66, VVD, PVV, CDA and JA21.}
\label{tab:party_within_between}
\begin{tabular}{@{}lll@{}}
\toprule
 & All parties & Six largest parties \\
\midrule
Within parties & 0.44 (0.03), $p < .001$ & 0.46 (0.03), $p < .001$ \\
Between parties & 0.09 (0.21), $p = .683$ & $-$0.68 (0.14), $p = .004$ \\
Between $-$ within & $-$0.36 (0.21), $p = .105$ & $-$1.14 (0.15), $p < .001$ \\
\addlinespace
Variance between parties (intercept) & 0.159 & 0.007 \\
Residual variance & 1.92 & 1.89 \\
Respondents & 1{,}153 & 923 \\
Parties & 14 & 6 \\
\bottomrule
\end{tabular}
\end{table}
